\section{Scalar compressions and coefficient centering}
\label{u:sec:centering}

Continue with a simple, unital, infinite-dimensional $A$ with strict
comparison and $T(A)\ne\varnothing$, and with $D$ and $\Lambda$ as in
Section~\ref{u:sec:preliminaries}.  We use the projection supply from
Proposition~\ref{u:prop:projection-supply} to partition most of a
projection into blocks on which prescribed coefficients have nearly
scalar compressions.  Trace-zero symmetries in those blocks will then
be centered against the coefficients.  All estimates are uniform over
$\Lambda$.  Traces on spectral projections in $D^{**}$ mean the normal
extensions of the traces on $D$; faithfulness of these extensions is
not required.

\begin{lemma}[Finite scalar compressions]
\label{u:lem:scalar-compressions}
Let $q\in D$ be a projection, let $x_1,\ldots,x_n\in D$ be selfadjoint
contractions, and let $\varepsilon,\delta>0$.  There are finitely many
pairwise orthogonal projections $p_1,\ldots,p_L\le q$ and numbers
$\mu_{s,j}\in[-1,1]$ such that
\begin{equation}
\label{u:eq:scalar-compressions}
 \begin{split}
 \norm{p_sx_jp_s-\mu_{s,j}p_s}&\le\varepsilon
       \quad(1\le j\le n,\ 1\le s\le L),\\
 \tau\!\left(q-\sum_{s=1}^L p_s\right)&\le\delta
       \quad(\tau\in\Lambda).
 \end{split}
\end{equation}
The family may be empty.  No commutation of the $p_s$ with the $x_j$ is
asserted.
\end{lemma}

\begin{proof}
It suffices to treat $0<\varepsilon,\delta<1$.  We first consider a
single selfadjoint contraction $x$.  Choose continuous functions
$f_1,\ldots,f_m:[-1,1]\to[0,1]$ and $\lambda_1,\ldots,\lambda_m\in[-1,1]$
such that
\[
 \sum_{j=1}^m f_j=1,
 \qquad
 \supp(f_j)\subseteq
 \{t\in[-1,1]:|t-\lambda_j|\le\varepsilon\}.
\]
For example, piecewise linear hat functions on a sufficiently fine
finite subdivision give such a partition of unity.

We will repeatedly apply these functions inside the current residual
corner and remove a projection supported where the corresponding
function is nonzero.  At any trace where a run starts with residual
trace greater than $\delta$, a full pass through all the functions will
remove a fixed positive trace mass.  To prove this uniformly, we first
estimate how the traces of the functional-calculus values change when
the corner shrinks.

If $v\in D$ is a projection, write $f^v(vxv)$ for functional calculus in the corner
$vDv$, whose unit is $v$; the value in the zero corner is defined to be
zero.  In particular, for a polynomial $P(t)=\sum_{k=0}^d a_kt^k$,
\begin{equation}
\label{u:eq:corner-polynomial}
 P^v(vxv)=a_0v+\sum_{k=1}^d a_k(vxv)^k.
\end{equation}

We give an explicit uniform continuity estimate for this changing
corner functional calculus.  Suppose $v\le u$ are projections and fix
$\tau\in\Lambda$.  Put $t=\tau(u-v)$.  Since $u-v$ is a projection,
\[
 \norm{u-v}_{2,\tau}=\sqrt t,
 \qquad
 \norm{uxu-vxv}_{2,\tau}\le2\sqrt t.
\]
The second inequality follows by writing the difference as
$(u-v)xu+vx(u-v)$ and using the tracial $2$-norm multiplication
inequalities.  The usual telescoping identity, with all powers of
$uxu$ and $vxv$ contractive, gives
\[
 \norm{(uxu)^k-(vxv)^k}_{2,\tau}\le 2k\sqrt t
 \qquad(k\ge1).
\]
Consequently, by $|\tau(a)|\le\norm{a}_{2,\tau}$ and
Equation~\eqref{u:eq:corner-polynomial},
\[
 \left|\tau\bigl(P^u(uxu)\bigr)
          -\tau\bigl(P^v(vxv)\bigr)\right|
 \le C_P\sqrt t,
 \qquad
 C_P=|a_0|+2\sum_{k=1}^d k|a_k|.
\]
The term $|a_0|$ records the change of corner units.  If
$f\in C([-1,1])$ and $\norm{f-P}_\infty\le a$, functional calculus in
each corner and $\tau(u),\tau(v)\le1$ yield
\begin{equation}
\label{u:eq:moving-corner-modulus}
 \left|\tau\bigl(f^u(uxu)\bigr)
          -\tau\bigl(f^v(vxv)\bigr)\right|
 \le 2a+C_P\sqrt{\tau(u-v)}.
\end{equation}
This estimate does not depend on $x$, the two projections, or the
chosen trace.

For each $j$ choose a polynomial $P_j$ with
$\norm{f_j-P_j}_\infty\le\delta/(16m)$, and set
\[
 C=\max\{1,C_{P_1},\ldots,C_{P_m}\},
 \qquad
 \alpha=\frac{\delta}{4m},
 \qquad
 \kappa=\min\left\{\frac{\delta}{4m},
                  \left(\frac{\delta}{8mC}\right)^2\right\}>0.
\]
Thus whenever $v\le u$ and $\tau(u-v)<\kappa$, Equation~
\eqref{u:eq:moving-corner-modulus} gives, for every $j$,
\begin{equation}
\label{u:eq:moving-corner-bound}
 \left|\tau\bigl(f_j^u(uxu)\bigr)
          -\tau\bigl(f_j^v(vxv)\bigr)\right|
 \le\frac{\delta}{4m}.
\end{equation}

Starting with an arbitrary residual projection $u$, perform one run
through the functions as follows.  Put $v_0=u$ and, successively, let
\[
 y_j=f_j^{v_{j-1}}(v_{j-1}xv_{j-1})\in v_{j-1}Dv_{j-1}
 \qquad(1\le j\le m).
\]
Apply Proposition~\ref{u:prop:projection-supply} to $y_j$ with cutoff
$\alpha$ to obtain a projection $a_j\in D$ satisfying
\[
 a_j\le\mathbf 1_{[\alpha,1]}(y_j)
       \le\supp(y_j)\le v_{j-1},
 \qquad
 \tau(a_j)=\tau((y_j-\alpha)_+)
           \ge\tau(y_j)-\alpha
 \quad(\tau\in\Lambda),
\]
and put $v_j=v_{j-1}-a_j$.  The trace estimate uses the stronger scalar
bound $0\le y_j-(y_j-\alpha)_+\le\alpha v_{j-1}$.  The projections
removed during the run are mutually orthogonal.  In the bidual corner,
$\supp(y_j)$ lies under the spectral projection of $v_{j-1}xv_{j-1}$
for $[\lambda_j-\varepsilon,\lambda_j+\varepsilon]\cap[-1,1]$.
Compressing the corresponding spectral order bounds by $a_j$ gives
\[
 -\varepsilon a_j\le a_jxa_j-\lambda_ja_j
                    \le\varepsilon a_j,
 \qquad
 \norm{a_jxa_j-\lambda_ja_j}\le\varepsilon.
\]
These are norm estimates in $D$ and do not require $a_j$ to commute
with $x$.

For every run and every $\tau\in\Lambda$ we claim that
\begin{equation}
\label{u:eq:scalar-run-progress}
 \tau(u)>\delta\quad\Longrightarrow\quad
 \tau(u-v_m)\ge\kappa.
\end{equation}
Indeed, suppose $\tau(u)>\delta$ but $\tau(u-v_m)<\kappa$.
The partition identity in the initial corner implies
\[
 \sum_{j=1}^m\tau\bigl(f_j^u(uxu)\bigr)=\tau(u)>\delta,
\]
so one index $j$ has $\tau(f_j^u(uxu))>\delta/m$.  Because
$v_{j-1}\le u$ and $\tau(u-v_{j-1})<\kappa$,
Equation~\eqref{u:eq:moving-corner-bound} applies at that step.  Therefore
\[
 \tau(a_j)\ge\tau(y_j)-\alpha
   >\frac{\delta}{m}-\frac{\delta}{4m}-\frac{\delta}{4m}
   =\frac{\delta}{2m}>\kappa,
\]
contradicting $a_j\le u-v_m$.  The index used to prove progress may
depend on the trace, but every run processes all $m$ indices and the
constant $\kappa$ is fixed beforehand.

Now start at $q$ and perform $K=\lfloor1/\kappa\rfloor+1$ runs,
each on the preceding residual projection.  Collect all removed
projections as $p_1,\ldots,p_L$, discarding zero projections if desired,
and let $w=q-\sum_s p_s$ be the final residual.  If some
$\tau\in\Lambda$ satisfied $\tau(w)>\delta$, the residual at the start
of every run would have trace greater than $\delta$ at this same
$\tau$.  Equation~\eqref{u:eq:scalar-run-progress} would imply
\[
 \tau(q-w)\ge K\kappa>1,
\]
which is impossible.  Thus $\tau(w)\le\delta$ simultaneously for all
$\tau\in\Lambda$.  This is a fixed finite construction; it uses no
compactness of $\Lambda$.  Each removed projection retains its scalar
compression estimate, proving the one-coefficient case.

For the finite family, begin with the single block $q$.  At stage
$j=1,\ldots,n$, suppose there are $L_{j-1}$ current blocks.  Apply the
one-coefficient case for $x_j$ inside each of them, with compression
tolerance $\varepsilon$ and trace-loss tolerance
\[
 \frac{\delta}{n\max\{1,L_{j-1}\}}.
\]
Replace the old blocks by their new subprojections.  The total trace
lost at this stage is at most $\delta/n$, at every trace.  All prior
compression bounds survive: if $p\le p_0$ and
$\norm{p_0x_kp_0-\mu p_0}\le\varepsilon$, then
\[
 px_kp-\mu p=p(p_0x_kp_0-\mu p_0)p,
 \qquad \norm{px_kp-\mu p}\le\varepsilon.
\]
There are finitely many blocks at every stage.  Summing the trace
losses over the $n$ stages proves Equation~\eqref{u:eq:scalar-compressions}.
For an empty coefficient family, simply take the block $q$.
\end{proof}

\begin{lemma}[Symmetries centered against finitely many coefficients]
\label{u:lem:coefficient-symmetry}
For a projection $q\in D$, a finite set $H\subset D$, and $\zeta>0$,
there is $r=r^*\in qDq$ such that
\begin{equation}
\label{u:eq:coefficient-symmetry}
 r^2=q,
 \qquad
 \sup_{\tau\in\Lambda}|\tau(br)|<\zeta
 \quad(b\in H).
\end{equation}
\end{lemma}

\begin{proof}
The assertion is immediate if $q=0$ or $H$ is empty.  Otherwise put
$M=\max\{1,\max_{b\in H}\norm{b}\}$ and choose
$0<\theta\le\min\{1/4,\zeta/(8M)\}$.  Apply
Lemma~\ref{u:lem:scalar-compressions}, with both tolerances $\theta$,
to the real and imaginary parts of every $b/M$, $b\in H$.
It gives orthogonal projections $p_s\le q$, a residual
$w=q-\sum_s p_s$, and complex scalars $\nu_{s,b}$ such that
\[
 \norm{p_sbp_s-\nu_{s,b}p_s}\le2M\theta,
 \qquad \tau(w)\le\theta\quad(\tau\in\Lambda).
\]
By Corollary~\ref{u:cor:corner-symmetry}, choose
$r_s=r_s^*\in p_sDp_s$ with $r_s^2=p_s$ and
$\tau(r_s)=0$ for every $\tau\in\Lambda$.  Set
\[
 r=w+\sum_s r_s.
\]
Orthogonality of the supporting projections gives $r=r^*\in qDq$ and
$r^2=w+\sum_s p_s=q$.  For a trace $\tau$, its restriction to
$p_sDp_s$ has norm $\tau(p_s)$.  Traciality therefore gives
\[
 \begin{split}
 |\tau(br_s)|
 &=\left|\tau\bigl((p_sbp_s-\nu_{s,b}p_s)r_s\bigr)\right|
 \le2M\theta\,\tau(p_s),\\
 |\tau(bw)|&=|\tau(wbw)|\le M\tau(w).
 \end{split}
\]
Summing these estimates, using $\sum_s\tau(p_s)\le1$, yields
\[
 |\tau(br)|\le2M\theta\sum_s\tau(p_s)+M\theta
             \le3M\theta\le\frac{3\zeta}{8}<\zeta.
\]
The error is weighted by the traces of the blocks, so their number
does not enter the bound.
\end{proof}

\begin{proposition}[Centered symmetries in finite-dimensional commutants]
\label{u:prop:centered-symmetry}
Let $E\subset D$ be a finite-dimensional $C^*$-subalgebra, possibly
nonunital, and let $S\subset D$ be a unital norm-separable
$C^*$-subalgebra.  There is a selfadjoint unitary $r\in D\cap E'$ such
that
\begin{equation}
\label{u:eq:centered-symmetry}
 \tau(sr)=0\qquad(s\in S,\ \tau\in\Lambda).
\end{equation}
\end{proposition}

\begin{proof}
Adjoin $1_D$ to $E$.  This preserves finite dimensionality: if $e$ is
the unit of $E$, it adds only the scalar summand $\C(1_D-e)$ when
$e\ne1_D$ (and if $E=0$, take $e=0$).  It therefore suffices to work
with the resulting unital finite-dimensional algebra.  Choose matrix
units
\[
 \{e^{(\ell)}_{ij}:1\le i,j\le n_\ell\},\qquad
 1\le\ell\le a,
 \qquad
 \sum_{\ell=1}^a\sum_{i=1}^{n_\ell}e^{(\ell)}_{ii}=1_D.
\]

We first impose finitely many coefficient tests.  Fix a finite
$H\subset S$ and $\xi>0$, and write $N_0=\sum_\ell n_\ell$.
For each block apply Lemma~\ref{u:lem:coefficient-symmetry} in the corner
$e^{(\ell)}_{11}De^{(\ell)}_{11}$ to the finite family
\[
 H_\ell=
 \{e^{(\ell)}_{1i}s e^{(\ell)}_{i1}:
                    s\in H,\ 1\le i\le n_\ell\},
\]
with tolerance $\xi/N_0$.  Denote the resulting corner symmetry by
$u_\ell$, and put
\[
 R=\sum_{\ell=1}^a\sum_{i=1}^{n_\ell}
       e^{(\ell)}_{i1}u_\ell e^{(\ell)}_{1i}.
\]
Each summand is selfadjoint, and the matrix-unit relations give
\[
 R^2=\sum_{\ell,i}e^{(\ell)}_{ii}=1_D,
 \qquad
 Re^{(\ell)}_{jk}
   =e^{(\ell)}_{j1}u_\ell e^{(\ell)}_{1k}
   =e^{(\ell)}_{jk}R.
\]
Terms from distinct matrix blocks have orthogonal supports.  Thus $R$
is a selfadjoint unitary commuting with $E$.  For every $s\in H$ and
$\tau\in\Lambda$, cyclicity gives
\[
 \tau(sR)=\sum_{\ell=1}^a\sum_{i=1}^{n_\ell}
       \tau\bigl(e^{(\ell)}_{1i}s e^{(\ell)}_{i1}u_\ell\bigr),
 \qquad
 |\tau(sR)|<N_0\frac{\xi}{N_0}=\xi.
\]
The complementary scalar block in the nonunital case is included in
these formulas, so the square of $R$ is the full unit of $D$.

Choose a norm-dense sequence $(s_k)_{k\ge1}$ in the unit ball of $S$.
Apply Lemma~\ref{u:lem:tracial-diagonal} to one variable in the
selfadjoint unit ball of $D$, with the tests
\[
 \norm{r^2-1_D}_{2,\Lambda}=0,
 \qquad
 \norm{[r,e^{(\ell)}_{ij}]}_{2,\Lambda}=0,
 \qquad
 \sup_{\tau\in\Lambda}|\tau(s_kr)|=0\quad(k\ge1).
\]
Every finite list is satisfiable to arbitrary precision by the
unitaries $R$ just constructed; the square and commutation relations
already hold exactly for those trial solutions.  The diagonal lemma
provides $r=r^*\in D$ satisfying every test.  Since the uniform
tracial $2$-norm detects zero in $D$, we have $r^2=1_D$ and
$[r,e^{(\ell)}_{ij}]=0$ in $D$.  Thus $r$ is a selfadjoint unitary in
$E'$.  Finally,
\[
 |\tau((s-s_k)r)|\le\norm{s-s_k}\norm{r}
\]
extends the trace identities from the dense sequence to the unit ball
of $S$, and scaling gives Equation~\eqref{u:eq:centered-symmetry} for
all $s\in S$.
\end{proof}

The proposition supplies exact commutation with $E$ and exact
coefficient centering over $S$.  It does not assert commutation with
$S$.  In the next section, this distinction permits successive choices
of symmetries centered against the algebra generated by their
predecessors.
